\scanpage{40}\mspage{$\underline{\overline{\mathrm{XIII}}}$}
\uline{Demonstration} : Pour voir que $\operatorname{int}\Lambda^2\subset\deleted{\unsure{\operatorname{ext}}}\operatorname{int}\Lambda^1$ \inserted{on} observe d'abord que :
\[
\Lambda^1-\{s,s'\}\subset\operatorname{ext}\Lambda^2\quad\text{donc}\quad\deleted{\unsure{\operatorname{ext}}}\operatorname{int}\Lambda^2\cap(\Lambda^1-\{s,s'\})=\varnothing
\]
, et par conséquent :
\[
\operatorname{int}\Lambda^2\subset\operatorname{int}\Lambda^1\quad\text{ou bien}\quad\operatorname{int}\Lambda^2\subset\operatorname{ext}\Lambda^1\quad(\text{car }\operatorname{int}\Lambda^2\subset\mathbb R^2-\Lambda^1).
\]
Mais $\operatorname{int.}\Lambda^2\cap\operatorname{int}\Lambda^1\ne\varnothing$. En effet soit $x\in{]s,s'[}$ et $D_x$ un disque de centre $x$ qui vérifie b) de la def.1 par rapport à $\Lambda^1$, $\Lambda^2$ et $\partial E(\Gamma)$, le lemme $\underline{\overline{\mathrm{XIII}}}$ montre que $D_x\cap\operatorname{ext}\partial E(\Gamma)=D_x\cap\operatorname{ext}\Lambda^1=D_x\cap\operatorname{ext}\Lambda^1$ donc $D_x\cap\deleted{\unsure{\operatorname{ext}}}\operatorname{int}\Lambda^1=D_x\cap\deleted{\unsure{\operatorname{ext}}}\operatorname{int}\Lambda^2$, et par conséquent $\operatorname{int}\Lambda^1\cap\operatorname{int}\Lambda^2\ne\varnothing$, c.q.f.d.

On veut prouver que $\Gamma$ est un $J$-polygone par récurrence sur $\operatorname{card}S(\Gamma)-2$.

Si $\operatorname{card}S(\Gamma)-2=1\quad\Rightarrow\quad\Gamma={}$ au bord d'un triangle\inserted{et le \unsure{theorēme|theorème} 1 devient donc trivial}, donc il suffit de prouver que si le \unsure{théorĕme|théorème} 1 est vrai pour tout $\Gamma$ avec $\operatorname{card}S(\Gamma)-2<n$ alors il est aussi vrai pour tout $\Gamma$ avec $\operatorname{card}S(\Gamma)-2=n$.

Il est clair que $\operatorname{card}S(\Lambda^i)<\operatorname{card}S(\Gamma)$ \inserted{$i=1,2$}, autrement $\Lambda^i=\Gamma$ ce qui est absurd puisque $ {]s,s'[}\cap\Gamma=\varnothing$. Les lemmes $\underline{\overline{\mathrm{XII}}}$ et $\underline{\overline{\mathrm{XIV}}}$ \inserted{avec le lemme $\underline{\overline{\mathrm{XV}}}$} nous permettent de supposer sans perte de generalité que : a) $\operatorname{int}\Lambda^2\subset\operatorname{int}\Lambda^1$ ;

b) $\Lambda^2-\{s,s'\}=\Gamma^2_{s,s'}-\{s,s'\}\subset\operatorname{int}\Lambda^1$.

Posons donc :
\[
\begin{aligned}
V&=\operatorname{ext}\Lambda^1\cup\operatorname{int}\Lambda^2\cup{]s,s'[}\quad\text{et}\\
U&=\operatorname{int}\Lambda^1-(\operatorname{int}\Lambda^2\cup(\Gamma^2_{ss'}\cap\operatorname{int}\Lambda^2)).
\end{aligned}
\]
Par \unsure{hypothēse|hypothèse} inductive $\operatorname{ext}\Lambda^1$ est non-relativement-compact donc $V$ non plus et $U\subset\operatorname{int}\Lambda^1$ est relativement compact puisque $\operatorname{int}\Lambda^1$ est relativement compact par hypothēse de récurrence.

\figure{\makebox[\linewidth][r]{\begin{tikzpicture}[x=1cm,y=1cm]
\fill (0,0) -- (.11,.015) -- (.15,-.06) -- cycle;
\draw (.05,-.015) .. controls (.94,-.34) and (1.60,-.80) .. (3.8,-1.20);
\end{tikzpicture}}}
% archive sheet number: 39
